\begin{abstract}
Multi-turn dialogue with Large Language Models (LLMs) suffers from progressive attention decay toward system prompt goal tokens. The Goal Attention Ratio (GAR) declines by 70.8\% between early and mid-conversation turns, yet linear probes recover goal content from the residual stream with AUC\,=\,0.98, indicating that goal information persists but the attention mechanism fails to route it. We introduce ThreadGate, a per-layer sigmoid gate on the residual stream ($x_{l+1} = x_l + g_l \cdot f_l(x_l)$), trained jointly with next-token prediction (NTP) and an auxiliary GAR loss. A four-way ablation (ThreadGate, Uniform Gate, LoRA-only, Base) on Qwen-2.5-7B-Instruct reveals that training succeeds (teacher-forced GAR\,=\,0.49) but generalization fails: under autoregressive generation, ThreadGate's GAR drops to 0.016, a 30$\times$ gap. Gating contributes zero measurable benefit, with only 0.75 percentage points separating all fine-tuned variants on held-out data. A PCGrad experiment rules out gradient direction conflict as the cause. The failure traces to exposure bias: attention patterns optimized under teacher-forced prefixes do not transfer to model-generated text, and the NTP objective's token-level self-recovery does not extend to auxiliary attention losses. These findings provide the first empirical characterization of how exposure bias selectively destroys auxiliary objective generalization while leaving the primary objective intact.
\end{abstract}
